\documentclass[10pt,a4paper]{amsart}
\usepackage[margin=19mm]{geometry}
\usepackage{amsmath,amssymb,mathtools}
\usepackage[scr=rsfs,cal=euler]{mathalfa}
\usepackage{xfrac}
\usepackage[unicode,hidelinks,bookmarks=false]{hyperref}
\newcommand{\ie}{i.e.\ }
\newcommand{\cf}{cf.\ }
\newcommand{\resp}{respectively\ }
\newcommand{\Subsection}{\subsection}
\providecommand{\nobreakdash}{\nobreak\hbox{-}}
\setlength{\emergencystretch}{2em}
\multlinegap0pt
\usepackage{bm}
\usepackage{bbm}
\usepackage{dsfont}
\usepackage{xspace}
\usepackage{cancel}
\usepackage{mathtools}
\usepackage{amsthm}
\makeatletter
\@ifundefined{theorem}{\newtheorem{theorem}{Theorem}}{}
\@ifundefined{lemma}{\newtheorem{lemma}[theorem]{Lemma}}{}
\@ifundefined{proposition}{\newtheorem{proposition}[theorem]{Proposition}}{}
\makeatother
\newcommand{\C}{\mathbb C}
\newcommand{\Gr}{\mathrm{Gr}}
\newcommand{\Q}{\mathbb Q}
\newcommand{\V}{\mathbb V}
\newcommand{\W}{\mathbb W}
\newcommand{\Z}{\mathbb Z}
\newcommand{\an}{\mathrm{an}}
\newcommand{\bul}{\bullet}
\newcommand{\gl}{\mathbf{GL}}
\title[On the distribution of mixed Hodge loci]{On the distribution of mixed Hodge loci}
\author{Nazim Khelifa}
\date{Source: arXiv:2603.20272v2 (CC BY 4.0)}
\begin{document}
\maketitle

\noindent\emph{Source and licence.} On the distribution of mixed Hodge loci. Version arXiv:2603.20272v2 (CC BY 4.0), distributed under the Creative Commons Attribution 4.0 International licence, as recorded in the arXiv licence metadata for this exact version and shown on the abstract page https://arxiv.org/abs/2603.20272v2. Full rights notice: licenses/arxiv-2603.20272-CC-BY-4.0.txt.\par\smallskip

\noindent\textit{Editorial scope.} Counted unit: the Proposition whose source label is \texttt{None} in the article (source0.tex lines 1838--1840, source label \texttt{None}), delivered together with the article's own complete original proof of it (source0.tex lines 1841--1900, one \texttt{proof} environment of 60 lines). No mathematics is written, repaired or re-derived: statement and proof are the article's own text line for line. Required context retained uncounted: monoint (lines 1801--1806). Every definition, display and label of those lines is retained unchanged. Omitted from this package and not redistributed: 1--1800, 1807--1837, 1901--2916. Nothing inside a retained excerpt is omitted or altered: every retained body is delivered exactly as the article's lines are written, so no omission receipt of any kind accompanies this package. Citations: the retained excerpts carry no citation command at all, so no bibliography entry of the article is transcribed and no reference is redistributed. Reference adaptation: none was needed and none was made; every reference inside the delivered unit resolves inside it, so no unresolved reference or citation is produced. Printed slips kept literal: no printed word, symbol, hypothesis or number of the counted statement or of its proof was altered. Layout and class: the article's own class is \texttt{article} with packages including geometry, titlesec, amsthm, amssymb, amsmath, biblatex, bbm, babel, inputenc, fontenc, xcolor, hyperref, ifpdf, graphicx; the wrapper is the lane's pinned \texttt{amsart} editorial wrapper at 10pt on A4, which restates the theorem-like environments the retained text needs and adds the article's own macro definitions that the retained text needs (\texttt{\textbackslash C}, \texttt{\textbackslash Gr}, \texttt{\textbackslash Q}, \texttt{\textbackslash V}, \texttt{\textbackslash W}, \texttt{\textbackslash Z}, \texttt{\textbackslash an}, \texttt{\textbackslash bul}, \texttt{\textbackslash gl}), with extra wrapper packages (bm, bbm, dsfont, xspace, cancel, mathtools). The delivered numbering is the unit's own. Assets: no retained span contains a figure environment, a graphics-inclusion command, a PDF-page inclusion, a table or a TikZ picture; no image file is copied into this package, the package declares no asset dependency and no image file is redistributed with it. Licence: version arXiv:2603.20272v2 of this article is distributed under the Creative Commons Attribution 4.0 International licence, as recorded in the arXiv licence metadata for this exact version and shown on the abstract page https://arxiv.org/abs/2603.20272v2. A full rights notice is delivered at licenses/arxiv-2603.20272-CC-BY-4.0.txt. Characters: every retained excerpt is pure ASCII, so the delivered engine, XeLaTeX with its default Latin Modern text font, reports no missing character.
\begin{lemma}\label{monoint}
Let $s \in S^\an$. The monodromy representation $\rho : \pi_1(S^\an,s) \rightarrow \gl(\mathcal{V}_s)$ associated to $(\mathcal{V}, \nabla)$ at the basepoint $s$ is given, in the basis $(e_{-2}^{(1)}, \cdots, e_{-2}^{(n)}, e_0)$, by the formula:
\[
\forall \gamma \in \pi_1(S^\an,s), \rho(\gamma) = \left(\begin{array}{cccc}1  &  \cdots & 0 & (2\pi i)m_1(q,\gamma)  \\ \vdots & \ddots & \vdots & \vdots \\0 & 0 & 1 & (2\pi i)m_n(q, \gamma)  \\0 & 0 & 0 & 1\end{array}\right).
\]
\end{lemma}

\begin{proposition}
Assume that for any holomorphic map $i : \Delta \rightarrow \overline{S}^\an$ such that $i(\Delta^\times) \subset S^\an$, the composite $q \circ i$ is meromorphic at $0$. Then, the associated variation $(\V, \W_\bul, \mathcal{F}^\bul)$ is admissible.
\end{proposition}

\begin{proof}
Let $i : \Delta \rightarrow \overline{S}^\an$ be a holomorphic map such that $i(\Delta^\times) \subset S$. We need to prove that $i^\ast (\V, \W_\bul, \mathcal{F}^\bul)$ is admissible. We make the usual identification $\pi_1(\Delta^\times) \cong \Z$, and denote by $\rho_i$ the monodromy representation associated to $i ^\ast \V$. According to Lemma \ref{monoint}, there exists a $n$-uple of integers $(m_1,\cdots, m_n)$ such that in the basis $(e_{-2}^{(1)}, \cdots, e_{-2}^{(n)}, e_0)$, one has:
\[
T := \rho_i(1) =  \left(\begin{array}{cccc}1  &  \cdots & 0 & (2\pi i) m_1 \\ \vdots & \ddots & \vdots & \vdots \\0 & 0 & 1 & (2\pi i) m_n  \\0 & 0 & 0 & 1\end{array}\right).
\]
Clearly this shows that the monodromy is unipotent, and actually $(T-1)^2 = 0$ so that $N := \log(T) = T-1$. Besides, since the monodromy is trivial on the graded pieces, we are in the setting of the following:
\begin{lemma}
Let $V$ be a $\Q$-vector space endowed with a finite increasing filtration $W_\bul$ and a nilpotent endomorphism $N$ which preserves $W$. Assume that the induced endomorphism $\Gr^W(N) : \Gr^W(V) \rightarrow \Gr^W(V)$ is zero. Then $V$ admits a weight filtration for $N$ relative to $W_\bul$ if and only if for every $k \in \Z$
\[
N(W_k) \subset W_{k-2}.
\]
\end{lemma}
Let $z \in \Delta^\times$. We have $N(e_0) = \sum_{1 \leqslant k \leqslant n} m_k e_{-2}^{(k)} \in (\W_{-2})_z$. Therefore the lemma shows that $\V_z$ admits a weight filtration for $N$ relative to $(\W_\bul)_z$.

It remains to establish the extension property. We first recall the construction of Deligne's canonical extension. Consider the connection on $\mathcal{V}$ defined by
\[
D = d - \left(\begin{array}{cccc}0  &  \cdots & 0 & m_1\frac{dz}{z} - \frac{d q_1}{q_1}  \\ \vdots & \ddots & \vdots & \vdots \\0 & 0 & 0 & m_n\frac{dz}{z} - \frac{d q_n}{q_n}  \\0 & 0 & 0 & 0\end{array}\right).
\]
The pair $(\mathcal{V}, D)$ is easily seen to have trivial associated monodromy representation, and the $D$-horizontal sections of $\mathcal{V}$ are well defined on the whole punctured disk and give another trivialization of $\mathcal{V}$:
\[
\mathcal{V} = \Big(\bigoplus_{1 \leqslant k \leqslant n} \mathcal{O}_{\Delta^\times} \cdot e_{-2}^{(k)} \Big) \oplus \mathcal{O}_{\Delta^\times} \cdot \tilde{e_0}
\]
where $\tilde{e_0} = e_0 + \sum_{1 \leqslant k \leqslant n} \log(\frac{q_k(i(z))}{z^{m_k}})e_{-2}^{(k)}$.
The canonical extension of $\mathcal{V}$ to $\Delta$ is then defined as the extension as a constant bundle for the above trivialization:
\[
\mathcal{V} = \Big(\bigoplus_{1 \leqslant k \leqslant n} \mathcal{O}_{\Delta} \cdot e_{-2}^{(k)} \Big) \oplus \mathcal{O}_{\Delta} \cdot \tilde{e_0}
\]
By assumption, the map $^iq := q \circ (i \vert_{\Delta^\times})$ is meromorphic at $0$ i.e. its components $(^iq_k)_{1 \leqslant k \leqslant n}$ are meromorphic at $0$ as functions $\Delta^\times \rightarrow \C^\times$. For $k \in \{1,\cdots, n\}$, there is by definition a unique writing
\[
^iq_k(z) = z^{l_k} (a_k + f_k(z))
\]
where $l_k \in \Z$, $a_k \in \C^\times$ and $f_k : \Delta \rightarrow \C$ is a holomorphic map vanishing at $0$.
\begin{lemma}\label{egal}
For $k \in \{1,\cdots, n\}$, one has $l_k = m_k$.
\end{lemma}
\begin{proof}
Let $k \in \{1,\cdots, n\}$. Consider a continuous loop $\gamma : [0,1] \rightarrow \Delta^\times$ whose homotopy class is the generator $1$ of $\pi_1(\Delta^\times)$. By definition $m_k$ is the index of $^i\gamma_k := ^iq_k \circ \gamma$ around $0$. If $l_k = 0$, one therefore finds
\[
2 \pi i m_k = \int_{^i\gamma_k} \frac{dz}{z} = \int_\gamma \frac{q'(u)}{q(u)}du = \int_\gamma \frac{f'_k(u)}{a_k + f_k(u)}du
\]
which is zero by the theorem of residues since the integrand of the last integral is holomorphic on the disk. Since we assumed $l_k = 0$, this proves that $m_k = l_k$. If $l_k \neq 0$, the same formula gives
\[
2 \pi i m_k = \int_{^i\gamma_k} \frac{dz}{z} = \int_\gamma \frac{q'(u)}{q(u)}du = \int_\gamma \frac{l_k du}{u} du + \int_\gamma \frac{f'_k(u)}{a_k + f_k(u)}du
\]
The second integral in the right hand side is zero as explained in the $l_k = 0$ case. The first integral in the right hand side is $2 \pi i l_k$ times the the index of $\gamma$ which is $1$ by definition. This finishes the proof.
\end{proof}
It follows from the above Lemma \ref{egal} that for $k \in \{1,\cdots, n\}$, one has for $z \in \Delta^\times$:
\[
\log(\frac{^iq_k(z)}{z^{m_k}}) = \log(a_k + f_k(z))
\]
which has a finite limit as $z \rightarrow 0$. Define the extension with the following formula, which makes sense by what precedes:
\[
\overline{\mathcal{F}^0} = \mathcal{O}_\Delta \cdot \Big( \tilde{e_0} - \sum_{1 \leqslant k \leqslant n} \log(\frac{q_k(i(z))}{z^{m_k}})e_{-2}^{(k)} \Big)
\]
This defines a locally free $\mathcal{O}_\Delta$-submodule of $\overline{\mathcal{V}}$, which is everywhere a supplement of 
\[
\overline{\mathcal{W}_{-2}} = \bigoplus_{1 \leqslant k \leqslant n} \mathcal{O}_{\Delta} \cdot e_{-2}^{(k)}.
\]
This is what we wanted to prove.
\end{proof}

\end{document}
